%%%%%%
%
% $Autor: Wings $
% $Datum: 2020-01-18 11:15:45Z $
% $Pfad: WuSt/Skript/Produktspezifikation/powerpoint/ImageProcessing.tex $
% $Version: 4620 $
%
%%%%%%



\chapter{Optimierungen}

In diesem Kapitel werden Möglichkeiten zur Optimierung eines neuronalen Netzes aufgezeigt
und anhand des bestehenden Netzes angewendet. Eine Optimierung lässt sich
hinsichtlich der Zeit beim Trainingsvorgang oder der Genauigkeit der Erkennung einer
Kategorie durchführen. Da es sich bei dem Trainingsmodell um ein einfaches neuronales
Netz handelt, ist der benötigte Rechenaufwand verhältnismäßig gering und der zeitliche
Aufwand überschaubar. Die folgenden Optimierungen befassen sich aus diesem Grund
verstärkt mit der Erhöhung und Festigung der Genauigkeit, welche im nachfolgenden
auch als Korrektklassifizierungsrate bezeichnet wird.

\section{Ausgangsbasis der künstlichen neuronalen Netze}

Das neuronale Netz zur Erkennung von Entnahmefehlern und das Netz zur Erkennung des
Fehlers Überschäumer, sind identisch aufgebaut. An dem Computer an der TS-30 wurde
das neuronale Netz zur Erkennung von Entnahmefehlern bereits aktiv getestet. Die 
eingehenden Bilder werden beim Ausführen des Modells in einer bestimmten Ordnerstruktur
gespeichert. Zuerst wird automatisch ein Ordner mit aktuellem Datum erstellt, falls dieser
noch nicht vorhanden ist. Danach wird die linke und rechte Bildhälfte auf einen 
Entnahmefehler geprüft. Das gesamte Bild wird nach der Auswertung in den entsprechenden
Ordner der Kategorie good good, good bad, bad good oder bad bad einsortiert und 
gespeichert. Es wird das Entnahmemodell für die nachfolgenden Optimierungen verwendet,
da die Genauigkeit direkt an dem Computer, der mit der TS-30 verbunden ist, getestet
werden kann.

Das ursprüngliche neuronale Netz zur Erkennung von Entnahmefehlern besaß folgende
Parameter:

\begin{itemize}
  \item Bildauflösung der Eingangsdaten: 80x30
  \item Anzahl der Epochen: 40
  \item Anzahl der Schritte pro Epoche: 200
  \item Stapelgröße: 64
  \item Netzstruktur: Faltungsebenen 16-32-32
  \item Kernelgröße: 3x3
  \item Poolinggröße: 2x2
\end{itemize}
	
	
Das Modell besitzt eine theoretische Genauigkeit von über 95\%. An der TS-30 wurde dieses 
Modell mit neu erzeugten Bildern getestet. Beim Prüfen der kategorisierten Ordner
stellte sich heraus, dass viele Bilder falsch sortiert wurden, was auf eine starke Abweichung
der tatsächlichen Genauigkeit vom theoretischen Wert hinweist. Mit dem Programm Prediction 
aus Kapitel 7 wurde eine Vorhersage auf einen Datensatz von insgesamt 3.730
Bildern gemacht. Der Datensatz bestand aus neuen Bildern, die noch nicht zum Trainieren 
verwendet wurden. Daraus ergab sich, dass aus dem Ordner mit 3.552 Bildern der
Kategorie good 172 Bilder fälschlicherweise der Kategorie bad zugeordnet wurden. Aus
dem Ordner der Kategorie bad in dem sich 178 Bilder befanden, wurden 128 Bilder als
good erkannt. Das Modell hat eine starke Tendenz dazu, Bilder der Kategorie bad in die
Kategorie good einzuordnen. Die praktische Gesamtgenauigkeit des Modells lag bei 91,96
Prozent. Aus diesen Erkenntnissen wurde das neuronale Netz schrittweise optimiert.

\subsection{Hinzufügen von neuen Bildmerkmalen}

Das Programm der Vorhersage aus Kapitel 7 zeigt die Bilder an, die vom Modell in
die falsche Kategorie eingeordnet wurden. Nach Begutachtung dieser Bilder stellte sich
heraus, dass alle Einzelformen und bestimmte Doppelformen besonders häufig als Fehler
erkannt wurden, obwohl das Produkt nicht im Werkzeug verharrte und somit kein Fehler
auftrat. Beim Überprüfen der ursprünglichen Trainingsdatenbasis zeigte sich, dass die
falsch kategorisierten Formen nicht in den Trainingsdaten enthalten waren. Daraufhin
wurden die gesamten Trainingsdaten aktualisiert. In diesem Zug wurden Bilder, die falsch
gedreht oder sehr unscharf waren, entfernt. Anschließend wurde eine Testmenge von ca.
10\%.

Die theoretische Genauigkeit lag nach dem Training bei 93,33\%. Mit dem Testdatensatz
wurde dieses Ergebnis kontrolliert. Von den 713 Bildern aus der Kategorie good wurden
712 richtig erkannt, was einer Genauigkeit von 99,86 Prozent entsprach und somit über
der Genauigkeit des Netzes lag. Bei 111 Bildern der Kategorie bad wurden 34 Bilder	
nicht richtig erkannt, es ergab sich eine Genauigkeit von 69,37 Prozent. Im Vergleich
zum Ursprung hat die Genauigkeit sich in beiden Fällen verbessert. In den folgenden
Abschnitten werden die Veränderungen der Genauigkeit mit Grafiken dokumentiert.

\subsection{Veränderung der Auflösung}

Die erste Optimierung bestand darin die Auflösung der Trainingsdaten zu verändern.
In Abbildung 8.1 (a) wird die Klassifizierungsgenauigkeit bei einer Datenbasis mit 
einer Auflösung von 80x30 je Bildhälfte dargestellt. Die blaue Linie in den Grafiken zeigt
die Veränderung der Genauigkeit der Trainingsdaten in Abhängigkeit der Epochen an.
Die roten Punkte zeigen die Genauigkeit der Validierungsdaten nach jeder Epoche an.
Es ist zu erkennen, dass die Validierungsgenauigkeit starke Schwankungen aufweist und
deutlich höher ist als die Trainingsgenauigkeit. Die Vermutung liegt nahe, dass das Netz
übertrainiert ist und die Validierungsdaten auswendig gelernt wurden. Für mehr Detail
Abstraktion wurde die Auflösung auf 160x60 pro Bildhälfte verdoppelt. Das Training mit
doppelt so hoher Auflösung sorgte für eine nahezu Vervierfachung des Rechenvorgangs.
Dies ist auf den exponentiellen Anstieg der Pixelmenge zurückzuführen. Die theoretische
Genauigkeit des neuronalen Netzes erhöhte sich gering auf 93,73\%. Angewendet auf die
Testdatenmenge, wurde ein Bild in der Kategorie good und 30 Bilder in der Kategorie bad
falsch erkannt. Die praktische Genauigkeit hat sich auf 96,24\% verbessert. Die Abbildung
8.1 (b) mit der Auflösung 60x160 zeigen ebenfalls starke Schwankungen bei den 
Validierungsgenauigkeiten Es wurde ein weiterer Durchlauf gestartet, um zu überprüfen,
ob eine weitere Erhöhung der Trainingsparameter zu einer Verbesserung der Ergebnisse führt. 
Die Zeit hat sich im Vergleich zum Ursprungsmodell von 28 Sekunden pro Epoche auf 462
Sekunden erhöht. Aufgrund der exponentiellen Vergrößerung der Datenpunkte stieg die
Zeit zum Trainieren exponentiell an. Die Genauigkeit hat sich auf 50 Prozent 
verschlechtert. Eine mögliche Ursache dafür könnte sein, dass die Filtergröße der Faltungschicht mit
3x3 zu klein war, um lokale Muster zu erkennen. Eine weitere Vergrößerung der Bildgröße
würde zu einem exponentiellen Anstieg der Zeit führen, weshalb die Größe 60x160 als
Basis für die nächsten Optimierungen verwendet wird.

\subsection{Anpassung der Validierungsdaten}

Die Validierungsmenge im ursprünglichen Modell lag bei einer Anzahl von 300 Bildern.
Auf die Gesamtmenge der Trainingsdaten entspricht das einen Prozentsatz von 3,6 Prozent. 
Als weiterer Optimierungsschritt wurde die Validierungsmenge auf 20 Prozent der
Gesamttrainingsmenge hochgesetzt.

Im Vergleich von Abbildung 8.1(b) zu Abbildung 8.2 ist zu erkennen, dass die Streu-
ung der Genauigkeit der Validierungsdaten näher bei der Genauigkeit der Trainingsdaten
liegt. Die Erhöhung der Validierungsdaten führte offensichtlich zu einer Verringerung des
Overfitting. Gleichzeitig führte es zu einer leichten Verlängerung der Epochen-Zeit, da
die Samples am Ende einer Epoche zur Überprüfung der Ergebnisse eingespielt werden.
Angewendet auf den Testdatensatz, wurden alle Bilder aus der Kategorie good richtig
erkannt und 24 Bilder aus der Kategorie bad falsch erkannt. Hierdurch ergibt sich eine
Gesamtgenauigkeit des Modells von 97\%.

\subsection{Veränderung der Epochen und Durchgänge}

Nach den ersten Optimierungsschritten wurde geprüft, ob ein Training mit weniger 
Epochen und weniger Schritten pro Epoche zu einem Ergebnis mit gleicher Genauigkeit führt.
Zudem könnte eine Reduzierung der Trainingszeit realisiert werden, wodurch die 
Aktualisierungen des Netzes schneller durchgeführt werden kann.

Es wurden verschiedene Varianten getestet. Zunächst wurde die Anzahl der Schritte auf
100 reduziert und die Anzahl der Epochen blieb unverändert bei 40. Die theoretische 
Genauigkeit der Trainings- und Validierungsdaten war fast identisch. Angewendet auf die
Testmenge wurde aus der Kategorie good ein Bild und aus der Kategorie bad wurden 26
Bilder falsch erkannt, was eine Erhöhung der falsch eingeordneten Bilder um 3 Stück 
ergab. Die Anzahl der Schritte wurde aus diesem Grund wieder auf 200 Schritte pro Epoche
erhöht.

Im zweiten Schritt wurde die Anzahl der Epochen auf 25 reduziert, was zu einer zeitlichen
Verkürzung des Trainings führte. Die theoretische Genauigkeit der Trainingsdaten 
reduzierte sich auf 89,19\% und die Validierungs-Genauigkeit auf 91,80\%. Die beiden Grafiken
befinden sich in Abbildung 8.3. Angewendet auf die Testmenge wurde ein Bild flasch in
die Kategorie good und 53 Bilder falsch in die Kategorie bad einsortiert. Da die 
Reduzierung der Epochen mit einer Verschlechterung der Genauigkeit einherging wurden keine
weiteren Tests mit reduzierten Epochen gemacht. Es wurde weiterhin mit 200 Schritten
pro Epoche und 40 Epochen gearbeitet.

\subsection{Daten künstlich Vervielfältigen}

Ein künstliches Vervielfältigen von Daten ist besonders empfehlenswert, wenn eine Klas-
sifikation stark unterrepräsentiert ist. Es wird dadurch ein künstlicher Ausgleich zwischen
den Kategorien erzeugt, um das neuronale Netz effizienter trainieren zu können. Dasselbe
Prinzip funktioniert auch mit der Reduzierung von Daten aus der überrepräsentierten
Klasse. Dabei werden Daten aus der Klasse zufällig aus dem Trainingsdatensatz ausge-
lassen und so eine Angleichung zwischen den Kategorien erzielt.

\subsubsection{Random-Oversampling}

In dem Ursprungsmodell wurde bereits das Random-Oversampling als Vervielfältigungsmethode
genutzt. Der entsprechende Programmabschnitt befinde sich im Kapitel 6.
Das Random-Oversampling verändert die Bilder nicht, sondern dupliziert Bilder aus der
unterrepräsentierten Kategorie bis ein Gleichgewicht zwischen den Kategorien herrscht.
Die Auswahl der vervielfältigten Bilder erfolgt dabei zufällig. \cite{Lemaitre:2017}

\subsubsection{SMOTE-Oversampling}
Zu weiteren Optimierungszwecken wurde das SMOTE-Oversampling als alternative Art
zum Random-Oversampling verwendet. Bei diesem Verfahren werden Bilder nicht eins
zu eins dupliziert, sondern durch Interpolation verändert. Die Technik beruht auf den
k-nächsten Nachbarn eines Bildes, basierend auf den euklidischen Abstand der Daten-
punkte. Zwischen einer betrachteten Instanz und den k nächsten Nachbarn existieren k
Differenzvektoren. Diese Differenzvektoren werden jeweils mit einer zufälligen Zahl > 0
bis < 1 multipliziert. Anschließend wird dieser neue Differenzvektor auf die 
Merkmalsvektoren der Instanz addiert. Der Vorteil ist, dass keine identischen Bilder 
für das Training verwendet werden. \cite{Chakrabarty:2015}

Nach Anpassung des Programms von Random-Oversampling auf SMOTE-Oversampling,
wurde das Netz neu trainiert. Die theoretische Genauigkeit des Netzes lag bei 95,89\% und
damit über der theoretischen Genauigkeit des Netzes mit dem Random-Oversampling.
Die Validierungsgenauigkeit lag bei 97,84\% was eine Verbesserung des Netzes erwarten
ließ. Angewendet auf die Testmenge wurden jedoch mehr Bilder falsch erkannt. Aus der
Kategorie good wurde ein Bild falsch erkannt und aus der Kategorie bad wurden 43 Bilder
falsch erkannt, entsprechend ergibt sich eine praktische Gesamtgenauigkeit von 94,7\%. Es
wurde daraufhin mit dem Random-Oversampling weitergearbeitet.

\subsection{Farbe}
Ein weiterer Versuch die Genauigkeit zu erhöhen lag darin, die Bilder der Datenbasis
als Farbbilder zum Trainieren zu verwenden. In dem Trainingsprogramm aus Kapitel
7 die Funktion GRAYSCALE entfernt und die Bildtiefe auf 3 hoch gesetzt. Nach dem
Training lag die Genauigkeit der Testdaten bei 91,51\% und die Validierungs-Genauigkeit
bei 90,42\%.

Angewendet auf die Testmenge wurden 29 Bilder in der Kategorie good und 54 Bilder
in der Kategorie bad nicht richtig zugeordnet. Sowohl die theoretische und praktische
Genauigkeit waren deutlich schlechter als bei dem neuronalen Netz, welches mit nur einer
Farbe trainiert wurde. Es ist zu vermuten, dass in den einzelnen Farbkanälen weniger
eindeutige Unterscheidungsmerkmale zu erkennen, da die Ursprungsbilder generell sehr
wenig Farbe enthalten. Daher wurde die Farbtiefe 1 als Basis beibehalten.

\subsection{Erweiterung einer Faltungsebene}
Der ursprüngliche Netzaufbau wurde zur weiteren Optimierung um eine Faltungsschicht
mit 32 Filtern ergänzt. Die theoretische Genauigkeit lag bei 93,44\%. Angewendet auf die
Testmenge wurden alle Bilder der Kategorie good richtig erkannt und 59 Bilder aus der
Kategorie bad falsch eingeordnet. Dies ergibt eine Gesamtgenauigkeit von 92,84\%. Da das
Hinzufügen einer weiteren Faltungsebene zu einer Verschlechterung geführt hat, wurde
dieser Optimierungsversuch verworfen.

\subsection{Veränderung der Filteranzahl}

Ein weiterer Optimierungsschritt bestand darin, die Filteranzahl der letzten Schicht zu
erhöhen. Sie wurde von 32 Filter auf 64 angepasst. Die theoretische Genauigkeit lag bei
95,3\%. Aus der Kategorie bad wurden 24 Bilder falsch zugeordnet, aus der anderen 
Kategorie wurden alle Bilder richtig erkannt. Die praktische Gesamtgenauigkeit lag bei 97,09%.
Die Erhöhung der Filter auf 64 brachte eine Verbesserung. Die Abbildung 8.7 zeigt den
Trainingsverlauf.

\subsection{Veränderung der Kernelgröße}

Als weiterer Optimierungsschritt wurde die Filtergröße geändert. Im ersten Durchgang
wurde sie auf 2x2 eingestellt. Ein kleinerer Filter kann auch komplexere Merkmale in 
einem Bild erkennen, jedoch sind die zu extrahierenden Merkmale sehr lokal. Grundlegende
Bildkomponenten könnten dadurch schlechter als Ganzes erkannt werden. Die theoretische 
Genauigkeit lag bei 95,84\%. Auf die Testmenge angewendet wurden alle Bilder der
Kategorie good richtig erkannt und 31 Bilder aus der Kategorie bad. Die praktische 
Genauigkeit lag bei 96,24\%. Im nächsten Schritt wurde die Kernelgröße auf 5x5 erhöht. Die
theoretische Genauigkeit lag bei 92,45\%. Angewendet auf die Testmenge ergaben sich 3
Fehler in der Kategorie good und 17 Fehler in der Kategorie bad. Die praktische 
Gesamtgenauigkeit des Netzes lag bei 97,94\%. Das neuronale Netz hatte in dieser 
Konstellation zwar die beste praktische Genauigkeit erzielt, jedoch hatte das Modell 
eine Tendenz entwickelt, Bilder ohne Entnahmefehler als Fehler zu erkennen. Für die praktische Anwendung
an der Maschine ist nachteilig, da der Produktionsprozess unnötig unterbrochen werden
könnte. Daher wurde die Filtergröße bei 3x3 belassen.

\section{Transfer Learning}

Die Parameter des optimierten neuronalen Netzes zur Erkennung von Entnahmefehlern
wurden auf das neuronale Netz zur Erkennung von Überschäumer übertragen.

Als erweiterte Funktion soll das neuronale Netz bei dem Fehler des
Überschäumers zudem erkennen können, ob es sich um einen Fehler vom Typ Deckel auf oder 
Deckel zu handelt (vergleiche Abbildung 3.2 aus Kapitel 3). Das ursprüngliche Modell zur 
Erkennung des Fehlers Deckel auf/ Deckel zu, hatte eine theoretische Genauigkeit von 
99,67\%. Das Modell wurde auf eine Testmenge aus 245 Bilder der Kategorie Deckel zu 
und 19 Bilder aus der Kategorie Deckel auf angewendet. 45 Bilder der Kategorie Deckel 
zu wurden falsch erkannt, was eine praktische Genauigkeit von 87,8\% ergab. 
Da für die Unterscheidung der Kategorien Deckel auf und Deckel zu nur sehr wenige 
Bilder zum Trainieren vorhanden waren, wurde eine Optimierung mit Hilfe des 
Transfer Learning durchgeführt. Als Basismodell wurde das neuronale Netz zur 
Erkennung von Überschäumen verwendet. Dieses besitzt eine sehr hohe
Ähnlichkeit zum Modell zur Erkennung des Deckelproblems, was
ein Vorteil beim Transfer Learning ist.

Der Programmaufbau entspricht dem bekannten Programmabschnitt zur Erstellung des
neuronalen Netzes aus Kapitel 7. Die wichtigsten Ergänzungen des Programms werden
nachfolgend beschrieben.

\begin{lstlisting}
import tensorflow as tf
from tensorflow.python.keras.models import Model , load_model
\end{lstlisting}

Eine wichtige Funktion, die in dem Programm aus Kapitel 7 ergänzt werden muss, ist
das load model. Dieses ermöglicht das Einspielen von schon vorhandenen Modellen in das
Programm.

\begin{lstlisting}
model=load_model ('Model_Ueberschaeumer.model')

model.summary ()
len( model.layers )
\end{lstlisting}

Mit der Funktion load model wird das neuronale Netz aufgerufen, welches sich in dem 
Verzeichnis unter dem Namen befindet. Anschließend werden mit dem Befehl model.summary
die Netzstruktur und die dazugehörigen Parameter angezeigt.

\begin{lstlisting}
for layer in model.layers [ : 1 2 ] :
    layer.trainable=False
for layer in model.layers [ 1 2 : ] :
    layer.trainable=True
\end{lstlisting}

Wenn das Modell geladen ist, kann mit den zwei for-Schleifen bestimmt werden, welche
Schichten des neuronalen Netzes neu trainiert werden. Abbildung 8.10 zeigt die 
Struktur des Modells zur Erkennung von Überschäumern und die Anzahl der trainierbaren
Parameter, die sich nach und nach erhöhen.

Die Abbildung 8.9 zeigt die theoretische Genauigkeit beim Trainieren
(Trainingsgenauigkeit) und die praktische Genauigkeit anwenden auf die Testmenge.
Die Tabelle 8.1 zeigt an, wie viele Bilder in der Testmenge der jeweiligen Kategorie 
von dem neuronalen Netz falsch zugeordnet wurden.

Bei 0 trainierbaren Parametern bleiben alle Gewichte und Schichten des Modells zur Er-
kennung von Überschäumern gleich. Das Netz hat daher keine Informationen über die
Bilder mit den Kategorien Deckel auf und Deckel zu und ordnet diese willkürlich den
Kategorien Überschäumer und kein Überschäumer zu. So lässt sich die theoretische 
Wahrscheinlichkeit von 47,73\% erklären.

Bei 33 trainierbaren Parametern ändert sich der Output-Layer, was bedeutet, dass die
Kategorien Überschäumer / kein Überschäumer zu Deckel auf / Deckel zu wechseln. Der
Rest des Modells wurde nicht angepasst. Der Einbruch der praktische Genauigkeit auf
7,2\% lässt sich dadurch erklären, dass alle Bilder der Testmenge Deckel auf/Deckel zu
gleichzeitig Überschäumer sind und vom Modell verstärkt in eine Kategorie sortiert werden.

Die beste Performance liefert das neuronale Netz bei 75.904 trainierbaren Parametern.
Dies ist dadurch zu erklären, dass dem Modell ab dieser Schicht örtliche Informationen
aus dem Training mitbekommt. Zudem sind die Filter in dieser Schicht auf die Erkennung
von Details ausgelegt.

Die weitere Verschlechterung der Genauigkeit mit zunehmenden trainierbaren Parameter
ist durch die steigende Abstraktion der Filter in Kombination mit einer geringen Anzahl
an Trainingsdaten zu erklären.
